\documentclass[9.5pt,a4paper,landscape]{article}

\usepackage[utf8]{inputenc}
\usepackage[T1]{fontenc}
\usepackage[french]{babel}
\usepackage[a4paper,landscape,margin=0.8cm]{geometry}
\usepackage{lmodern}
\usepackage{xcolor}
\usepackage{tcolorbox}
\usepackage{amsmath}

\pagestyle{empty}
\setlength{\parindent}{0pt}

\definecolor{bleu}{HTML}{245B78}
\definecolor{gris}{HTML}{F7F7F7}

\newtcolorbox{exo}[1]{
  colback=white,
  colframe=black!55,
  boxrule=0.45pt,
  arc=1.2mm,
  left=1.8mm,
  right=1.8mm,
  top=1.5mm,
  bottom=1.3mm,
  height=5.35cm,
  valign=top,
  title=\textbf{Exercice hors classe #1},
  coltitle=bleu,
  fonttitle=\small,
  colbacktitle=gris,
  titlerule=0.25pt
}

\begin{document}

\fontsize{9}{10.5}\selectfont

\begin{center}
{\large\bfseries\color{bleu}
Séquence 2 -- Théorème de Pythagore : exercices hors classe}\par

\vspace{0.1em}

{\small
À faire dans le cahier.
Ne faire que l'exercice indiqué par le professeur.
Calculatrice autorisée.}
\end{center}

\vspace{0.25em}

%==================================================
% LIGNE 1
%==================================================

\noindent
\begin{minipage}[t]{0.242\linewidth}
\begin{exo}{1}
Le triangle \(MNP\) est rectangle en \(N\).

\medskip
1. Quelle est son hypoténuse ?

\medskip
2. Écris l'égalité de Pythagore avec les lettres.
\end{exo}
\end{minipage}
\hfill
%
\begin{minipage}[t]{0.242\linewidth}
\begin{exo}{2}
Le triangle \(ABC\) est rectangle en \(A\).

\[
AB=6\text{ cm}
\qquad
AC=8\text{ cm}
\]

Un élève affirme que \(BC=7\) cm.

Peut-il avoir raison ? Justifie en une phrase.
\end{exo}
\end{minipage}
\hfill
%
\begin{minipage}[t]{0.242\linewidth}
\begin{exo}{3}
Le triangle \(RST\) est rectangle en \(S\).

\[
RS=6\text{ cm}
\qquad
ST=8\text{ cm}
\]

Calcule la longueur \(RT\).
\end{exo}
\end{minipage}
\hfill
%
\begin{minipage}[t]{0.242\linewidth}
\begin{exo}{4}
Le triangle \(DEF\) est rectangle en \(E\).

\[
DF=13\text{ cm}
\qquad
EF=5\text{ cm}
\]

Calcule la longueur \(DE\).
\end{exo}
\end{minipage}

\vspace{0.16cm}

%==================================================
% LIGNE 2
%==================================================

\noindent
\begin{minipage}[t]{0.242\linewidth}
\begin{exo}{5}
Le triangle \(GHI\) est rectangle en \(H\).

\[
GH=7{,}2\text{ cm}
\qquad
HI=9{,}6\text{ cm}
\]

Calcule la longueur \(GI\).
\end{exo}
\end{minipage}
\hfill
%
\begin{minipage}[t]{0.242\linewidth}
\begin{exo}{6}
Le triangle \(JKL\) est rectangle en \(K\).

\[
JL=7{,}3\text{ cm}
\qquad
KL=4{,}8\text{ cm}
\]

Calcule la longueur \(JK\).
\end{exo}
\end{minipage}
\hfill
%
\begin{minipage}[t]{0.242\linewidth}
\begin{exo}{7}
Complète :

\[
9^2=81
\qquad\text{donc}\qquad
\sqrt{81}=\ldots
\]

\[
12^2=144
\qquad\text{donc}\qquad
\sqrt{144}=\ldots
\]

\[
1{,}5^2=2{,}25
\qquad\text{donc}\qquad
\sqrt{2{,}25}=\ldots
\]
\end{exo}
\end{minipage}
\hfill
%
\begin{minipage}[t]{0.242\linewidth}
\begin{exo}{8}
Le triangle \(PQR\) est rectangle en \(Q\) et isocèle en \(Q\).

\[
PQ=QR=6\text{ cm}
\]

Donne la valeur exacte de \(PR\), puis un arrondi au dixième de cm.
\end{exo}
\end{minipage}

\vspace{0.16cm}

%==================================================
% LIGNE 3
%==================================================

\noindent
\begin{minipage}[t]{0.242\linewidth}
\begin{exo}{9}
Le triangle \(UVW\) est rectangle en \(V\).

\[
UW=10\text{ cm}
\qquad
VW=7\text{ cm}
\]

Donne la valeur exacte de \(UV\), puis un arrondi au dixième de cm.
\end{exo}
\end{minipage}
\hfill
%
\begin{minipage}[t]{0.242\linewidth}
\begin{exo}{10}
Complète les phrases.

\medskip

Le carré de \(7\) est \ldots

\medskip

La racine carrée de \(49\) est \ldots

\medskip

Le nombre \ldots{} a pour carré \(36\).

\medskip

Le nombre \ldots{} est la racine carrée de \(16\).
\end{exo}
\end{minipage}
\hfill
%
\begin{minipage}[t]{0.242\linewidth}
\begin{exo}{11}
Écris chaque phrase sous forme d'une égalité.

\medskip

1. \(9\) est la racine carrée de \(81\).

\medskip

2. \(49\) est le carré de \(7\).

\medskip

3. \(5\) est la racine carrée de \(25\).
\end{exo}
\end{minipage}
\hfill
%
\begin{minipage}[t]{0.242\linewidth}
\begin{exo}{12}
Le triangle \(IJK\) est rectangle en \(K\) avec

\[
IK=7\text{ cm}
\qquad
JK=2{,}4\text{ cm}.
\]

Calcule \(IJ\) en rédigeant une solution complète.
\end{exo}
\end{minipage}

%==================================================
% CORRIGÉ PROF (non imprimé)
%==================================================
%
% 1) Hypoténuse : [MP].
%    MP^2 = MN^2 + NP^2.
%
% 2) Non : l'hypoténuse BC doit être le plus long côté ;
%    or 7 < 8.
%
% 3) RT^2 = 6^2 + 8^2 = 100
%    donc RT = 10 cm.
%
% 4) DE^2 = 13^2 - 5^2 = 144
%    donc DE = 12 cm.
%
% 5) GI^2 = 7,2^2 + 9,6^2 = 144
%    donc GI = 12 cm.
%
% 6) JK^2 = 7,3^2 - 4,8^2 = 30,25
%    donc JK = 5,5 cm.
%
% 7) 9 ; 12 ; 1,5.
%
% 8) PR^2 = 6^2 + 6^2 = 72
%    donc PR = \sqrt{72} cm
%    et PR \approx 8,5 cm.
%
% 9) UV^2 = 10^2 - 7^2 = 51
%    donc UV = \sqrt{51} cm
%    et UV \approx 7,1 cm.
%
% 10) 49 ; 7 ; 6 ; 4.
%
% 11) \sqrt{81}=9
%     7^2=49
%     \sqrt{25}=5
%
% 12) Le triangle IJK est rectangle en K.
%     D'après le théorème de Pythagore :
%     IJ^2 = IK^2 + JK^2
%     IJ^2 = 7^2 + 2,4^2
%     IJ^2 = 54,76
%     IJ = \sqrt{54,76}
%     donc IJ = 7,4 cm.
%
\end{document}